\documentclass[12pt,a4paper]{article}
\usepackage[utf8]{inputenc}
\usepackage[indonesian]{babel}
\usepackage{amsmath}
\usepackage{amsfonts}
\usepackage{amssymb}
\usepackage{geometry}

\geometry{margin=2.5cm}

\title{Menghitung Luas Permukaan Benda Putar}
\author{Ahmad Fakhrul Bawani}
\date{}

\begin{document}

\maketitle

\section*{Soal}

Find the area of the surface generated by revolving the curve:
\begin{equation*}
  y = \frac{x^3}{9}, \quad \text{untuk } 0 \le x \le 2, \quad \text{mengelilingi sumbu-}x
\end{equation*}

\subsection*{1. Rumus Umum}
Rumus luas permukaan ($S$) benda putar yang dihasilkan oleh perputaran mengelilingi sumbu-$x$ dari $x = a$ sampai $x = b$ didefinisikan sebagai:
\begin{equation*}
  S = 2\pi \int_{a}^{b} y \cdot \sqrt{1 + \left(\frac{dy}{dx}\right)^2} \, dx
\end{equation*}

\subsection*{2. Menentukan Turunan Fungsi}
Cari turunan pertama fungsi $y = \frac{1}{9}x^3$ terhadap $x$:
\begin{equation*}
  \frac{dy}{dx} = \frac{1}{9}(3x^2) = \frac{x^2}{3}
\end{equation*}

Kuadratkan hasil turunan tersebut:
\begin{equation*}
  \left(\frac{dy}{dx}\right)^2 = \left(\frac{x^2}{3}\right)^2 = \frac{x^4}{9}
\end{equation*}

Masukkan ke dalam bentuk akar kuadrat pembentuk panjang busur:
\begin{equation*}
  \sqrt{1 + \left(\frac{dy}{dx}\right)^2} = \sqrt{1 + \frac{x^4}{9}} = \sqrt{\frac{9 + x^4}{9}} = \frac{1}{3}\sqrt{9 + x^4}
\end{equation*}

\subsection*{3. Solusi Integrasi}
Substitusikan komponen $y$ dan bentuk akar di atas ke dalam rumus luas permukaan dengan batas $a = 0$ dan $b = 2$:
\begin{align*}
  S &= 2\pi \int_{0}^{2} \left(\frac{x^3}{9}\right) \cdot \left(\frac{1}{3}\sqrt{9 + x^4}\right) \, dx \\
  S &= \frac{2\pi}{27} \int_{0}^{2} x^3 (9 + x^4)^{\frac{1}{2}} \, dx
\end{align*}

Gunakan metode substitusi untuk menyelesaikan integral tersebut:
\begin{align*}
  \text{Misalkan: } & u = 9 + x^4 \implies du = 4x^3 \, dx \implies x^3 \, dx = \frac{1}{4} \, du \\
  \text{Batas baru: } & \text{Untuk } x = 0 \implies u = 9 + 0^4 = 9 \\
  & \text{Untuk } x = 2 \implies u = 9 + 2^4 = 9 + 16 = 25
\end{align*}

Substitusikan variabel baru $u$ ke dalam integral:
\begin{align*}
  S &= \frac{2\pi}{27} \int_{9}^{25} u^{\frac{1}{2}} \cdot \left(\frac{1}{4} \, du\right) \\
  S &= \frac{2\pi}{108} \int_{9}^{25} u^{\frac{1}{2}} \, du \\
  S &= \frac{\pi}{54} \left[ \frac{2}{3}u^{\frac{3}{2}} \right]_{9}^{25} \\
  S &= \frac{\pi}{81} \left[ u\sqrt{u} \right]_{9}^{25} \\
  S &= \frac{\pi}{81} \left( \left[ 25\sqrt{25} \right] - \left[ 9\sqrt{9} \right] \right) \\
  S &= \frac{\pi}{81} \left( 25(5) - 9(3) \right) \\
  S &= \frac{\pi}{81} (125 - 27) \\
  S &= \frac{\pi}{81} (98) \\
  S &= \frac{98\pi}{81}
\end{align*}

Jadi, luas permukaan benda putar tersebut adalah \textbf{$\dfrac{98\pi}{81}$ satuan luas}.

\end{document}